\documentclass[10pt,a4paper]{amsart}
\usepackage[margin=19mm]{geometry}
\usepackage{amsmath,amssymb,mathtools}
\usepackage[scr=rsfs,cal=euler]{mathalfa}
\usepackage{xfrac}
\usepackage[unicode,hidelinks,bookmarks=false]{hyperref}
\newcommand{\ie}{i.e.\ }
\newcommand{\cf}{cf.\ }
\newcommand{\resp}{respectively\ }
\newcommand{\Subsection}{\subsection}
\providecommand{\nobreakdash}{\nobreak\hbox{-}}
\setlength{\emergencystretch}{2em}
\multlinegap0pt
\usepackage{amssymb}
\usepackage{mathtools}
\usepackage[shortlabels]{enumitem}
\usepackage{tikz}
\newtheorem{lemma}{Lemma}
\title{Residue Constraints in the Rank-Three Lifting Problem for Projective-Plane Incidence Matrices}
\author{Jaehwan Kim}
\date{Source: arXiv:2605.08090v1 (CC BY 4.0)}
\begin{document}
\maketitle

\noindent\emph{Source and licence.} Residue Constraints in the Rank-Three Lifting Problem for Projective-Plane Incidence Matrices. Version arXiv:2605.08090v1 (CC BY 4.0), distributed by arXiv under the Creative Commons Attribution 4.0 International licence, http://creativecommons.org/licenses/by/4.0/, as recorded in the arXiv licence metadata for this exact version and shown on the abstract page https://arxiv.org/abs/2605.08090v1. Full licence notice: \texttt{licenses/arxiv-2605.08090-CC-BY-4.0.txt}.\par\smallskip

\noindent\textit{Editorial scope.} Counted unit: the article's lemma lem:identity-count-detailed (source0.tex lines 3528--3536). The article's complete original proof of it is delivered with it (source0.tex lines 3538--3588, one \texttt{proof} environment of 51 lines). No mathematics is written, repaired or re-derived: the statement and the proof are the article's own lines, in the article's own notation, and no retained line is substituted or re-flowed.  No further source-local context is needed: every cross-reference of every retained line resolves inside the two delivered spans.  Nothing was omitted from any retained span, so the package carries no omission record.  No retained line contains a citation, so the wrapper carries no bibliography and no bibliography entry.  No retained line contains a figure environment, an image-inclusion command or drawing code, so no image is delivered and the package declares no asset dependency.  Editorial formatting only, in addition: an AMS \texttt{amsart} preamble that restates the article's own declarations and macros used by the retained lines, namely the environments \texttt{lemma} on separate counters, together with the standard packages those lines need.  The retained environments print in this document as Lemma 1 in order of appearance, whereas the article prints the same items with the numbering of its own full document.  The complete native source of the article as distributed in the arXiv e-print for this exact version is bundled unchanged as \texttt{sources/200132/source0.tex}.
\begin{lemma}[Detailed private-line count for identity-pattern minors]
\label{lem:identity-count-detailed}
Let $\Pi$ be a projective plane of order $q\ge 3$ with combinatorial incidence matrix $I(\Pi)$, and let $v:=q^2+q+1$.
The number of ordered $4\times4$ submatrices of $I(\Pi)$ equal to $I_4$ and obtained from ordered quadrangles with one private line through each vertex is at least
\[
v(v-1)q^2(q-1)^2(q-2)^4.
\tag{A.1}
\]
\end{lemma}

\begin{proof}
Choose an ordered quadrangle $(P_1,P_2,P_3,P_4)$ first.

There are $v$ choices for $P_1$ and $v-1$ choices for $P_2\neq P_1$.
The third point $P_3$ must avoid the line $\overline{P_1P_2}$.
Because every line has $q+1$ points, the number of available choices is
\[
v-(q+1)=q^2.
\]
Thus $(P_1,P_2,P_3)$ is already noncollinear.

For $P_4$, exclude the three lines
\[
\overline{P_1P_2},\qquad \overline{P_1P_3},\qquad \overline{P_2P_3}.
\]
Each contains $q+1$ points.
Since $P_1,P_2,P_3$ are noncollinear, those lines are pairwise distinct and meet pairwise exactly at $P_1,P_2,P_3$.
Hence their union has size
\[
3(q+1)-3=3q.
\]
So the number of admissible choices for $P_4$ is at least
\[
v-3q=(q^2+q+1)-3q=(q-1)^2.
\]
Any such $P_4$ yields an ordered quadrangle.

Now fix such a quadrangle.
For each $i\in\{1,2,3,4\}$, exactly $q+1$ lines pass through $P_i$.
Among them, precisely three also pass through one of the other vertices, namely
\[
\overline{P_iP_j}\qquad (j\neq i).
\]
These three lines are distinct because no three vertices are collinear.
Therefore exactly
\[
(q+1)-3=q-2
\]
lines through $P_i$ avoid the other three vertices.
Choosing one such line $L_i$ for each $i$ produces a $4\times4$ incidence block with
\[
P_i\in L_i,\qquad P_j\notin L_i\quad (j\neq i),
\]
hence the block is exactly $I_4$.

Multiplying the four point-selection factors and the four private-line factors gives
\[
v(v-1)\cdot q^2\cdot (q-1)^2\cdot (q-2)^4,
\]
which is \textup{(A.1)}.
\end{proof}

\end{document}
